\section{List coloring small graphs}
\label{small:section}

\begin{proposition}\label{prop:small-list}
There are an absolute positive integer $L$ and an absolute integer $t_0$
such that, for every integer $t\ge t_0$ and every graph $G$,
\[
 h(G)<t,\qquad |V(G)|\le t^{11/10}
 \quad\Longrightarrow\quad \ell(G)\le Lt.
\]
\end{proposition}

We will color the graph in stages, halving an upper bound on its fractional
chromatic number at each stage. For each sampled stable set, only a specified
subset of vertices is eligible to be removed. The next lemma allows
these eligibility restrictions to be fixed in advance. In the final
application they will come from one partition of the actual colors
appearing in the lists, chosen before the successive remaining graphs
are determined.

\begin{lemma}[Eligible thinning]\label{small:eligible-thinning}
There is an absolute positive integer $B$ such that the following holds
for all sufficiently large integers $t$. Let $n,r$ be real parameters with
\[
 1\le n\le t^{11/10},\qquad \frac{t}{\log t}\le r\le4t,
\]
and let $P$ be a graph with $|V(P)|\le n$, $h(P)<t$, and $f(P)\le r$.
Set
\[
 N=\left\lceil Br\log\bigl(2\max\{1,t/r\}\bigr)\right\rceil.
\]
Suppose $E_1,\ldots,E_N\subseteq V(P)$ contain every vertex at least
$N/2$ times. There are stable sets $I_1,\ldots,I_N$ of $P$, each of size
at most $\lceil n/r\rceil$, such that
\[
 f\left(P-\bigcup_{j=1}^N(I_j\cap E_j)\right)\le r/2.
\]
The cutoff for $t$ is uniform over $P,n,r$ and the eligibility sets.
\end{lemma}

We prove the lemma by contradiction. If every sampled remainder has
fractional chromatic number greater than $r/2$, we first find a fixed
highly connected subgraph $H$ where a remainder still has large expected
fractional chromatic number. Fresh independent remainders then contain
disjoint subsets, each carrying a clique model. The complement in $H$
of the union of the remainders remains highly connected. To assemble a
$K_t$ model, partition its $t$ labels into small blocks and use one of
these disjoint models for each
block and each pair of blocks. The complement connects all bags assigned
to the same label. Every pair of labels is represented together in one
model, so the resulting label sets are adjacent. The estimates below
provide the models and the connectivity for this construction.

\subsection{Eligible samples and uniform concentration}

We prove Lemma~\ref{small:eligible-thinning} with the following fixed
choice of $B$. Let $\gamma$ be the integer in
Theorem~\ref{thm:grouping}. Fix
\begin{equation}\label{small:constants}
 D_0=10^6\gamma,\qquad Q_0=401^2,
\end{equation}
and choose a positive integer $B$ such that
\begin{equation}\label{small:B-choice}
 B\ge16,\qquad
 2^{B/4}\ge
 \max\bigl\{2000(D_0+1),\,10^4Q_0^2\bigr\}.
\end{equation}
These constants remain fixed throughout this section.

Consider an integer $t$ and real parameters $n,r$ satisfying
\begin{equation}\label{small:range}
 1\le n\le t^{11/10},\qquad
 \frac{t}{\log t}\le r\le4t.
\end{equation}
Let $P$ be a graph on at most $n$ vertices with $h(P)<t$ and $f(P)\le r$.
Set
\begin{equation}\label{small:parameters}
 \begin{gathered}
 y=\max\{1,t/r\},\qquad x=\lfloor r/200\rfloor,
 \qquad q=\lceil t/x\rceil^2,\qquad k=\lceil D_0ty\rceil,\\
 N=\lceil Br\log(2y)\rceil,\qquad M=\lceil n/r\rceil.
 \end{gathered}
\end{equation}
Here $x$ bounds the size of a label block, $q$ bounds the number of
disjoint clique models, and $k$ is the connectivity required of $H$.
The parameters $N$ and $M$ are the number of trials and the maximum
size of one sampled stable set.
We always take $t$ above an absolute cutoff large enough that $r\ge400$.
In particular $x\ge1$. The following bounds hold uniformly in
\eqref{small:range}:
\begin{equation}\label{small:parameter-bounds}
 \begin{gathered}
 1\le y\le\log t,\qquad \frac{t}{y}\le r\le\frac{4t}{y},\qquad
 x\ge\frac r{400},\\
 q\le Q_0y^2,\qquad xq\le403ty,\qquad
 D_0ty\le k\le(D_0+1)ty.
 \end{gathered}
\end{equation}
Indeed, $t/x\le400t/r\le400y$, so
$\lceil t/x\rceil\le401y$. Also
\[
 x\lceil t/x\rceil^2
 \le\frac{t^2}{x}+2t+x
 \le400\frac{t^2}{r}+2t+\frac r{200}
 \le403ty.
\]
Our choice of $D_0$ consequently gives, for all sufficiently large $t$,
\begin{equation}\label{small:attachment-budget}
 k\ge r,\qquad 2xq<\frac{k}{64},\qquad
 \left\lfloor\frac{k}{32}\right\rfloor\ge\gamma(2xq).
\end{equation}

Fix sets $E_1,\ldots,E_N\subseteq V(P)$ such that each vertex belongs
to at least $N/2$ of them. These are the eligibility sets and are fixed
before sampling. Independently for each $j\in\{1,\ldots,N\}$, sample a
stable set $I_j$ from the distribution in Lemma~\ref{lem:bounded}, with
\[
 |I_j|\le M,\qquad \Pr(v\in I_j)=\frac1{2r}.
\]
The covered vertices at trial $j$ are $I_j\cap E_j$, and the remainder is
\begin{equation}\label{small:remainder}
 W=V(P)\setminus\bigcup_{j=1}^N(I_j\cap E_j).
\end{equation}
For every vertex $v$, independence of the trials and the eligibility
condition imply
\begin{equation}\label{small:survival}
 \Pr(v\in W)
 \le\left(1-\frac1{2r}\right)^{N/2}
 \le p:=\exp\left(-\frac{N}{4r}\right)
 \le(2y)^{-B/4}.
\end{equation}
In particular, \eqref{small:B-choice} and \eqref{small:parameter-bounds}
give
\begin{equation}\label{small:sparseness}
 2kp\le\frac r{1000},\qquad q^2p\le10^{-4}.
\end{equation}
For the first inequality, divide by $r$ and use
\[
 \frac{2kp}{r}
 \le2(D_0+1)y^2(2y)^{-B/4}
 \le2(D_0+1)2^{-B/4}.
\]
For the second, use
$q^2p\le Q_0^2y^4(2y)^{-B/4}\le Q_0^2 2^{-B/4}$.

\begin{lemma}\label{small:concentration}
Take at most $q$ independent copies of the experiment
\eqref{small:remainder}. Let $Z$ be any one of the following quantities,
where the vertex set $A\subseteq V(P)$ is fixed before sampling:
\[
 |A\cap W_i|,\qquad
 \left|A\cap\bigcup_iW_i\right|,\qquad
 f(P[A\cap W_i]).
\]
For each fixed $\delta>0$, and all sufficiently large $t$,
\begin{equation}\label{small:uniform-tail}
 \Pr\bigl(|Z-\mathbb EZ|\ge\delta r\bigr)
 \le2\exp(-t^{7/10}),
\end{equation}
uniformly over all the data in \eqref{small:range} and all eligibility
sets satisfying the stated condition.
\end{lemma}

\begin{proof}
Changing one sampled stable set changes $W_i$ on at most $2M$ vertices.
It changes a union of remainders on at most the same number of vertices.
Adding or deleting a vertex changes fractional chromatic number by at
most one, so all the quantities in the statement have bounded differences
at most $2M$ per trial.

For sufficiently large $t$, $r\log(2y)\ge1$, and hence
$N\le(B+1)r\log(2y)$. With
$A_0=4Q_0(B+1)$, we obtain
\[
 qN\le A_0ty\log(2y),\qquad M\le2t^{1/10}y.
\]
Lemma~\ref{lem:concentration} therefore gives
\begin{align*}
 \Pr\bigl(|Z-\mathbb EZ|\ge\delta r\bigr)
 &\le2\exp\left(-\frac{\delta^2r^2}{2(qN)(2M)^2}\right)\\
 &\le2\exp\left(
 -\frac{\delta^2t^{4/5}}{32A_0y^5\log(2y)}\right).
\end{align*}
Since $1\le y\le\log t$, the exponent in the last line has magnitude at
least $t^{7/10}$ for all sufficiently large $t$. The cutoff depends only
on $\delta$ and the fixed constants. Independence between different
vertices is not used.
\end{proof}

All asymptotic probability estimates below are uniform over these data.
Only finitely many fixed positive values of $\delta$ will be used in
Lemma~\ref{small:concentration}, so their cutoffs can be combined.

\subsection{Localizing a failure of thinning}

\begin{lemma}\label{small:fixed-region}
Suppose that every outcome of the experiment
\eqref{small:remainder} having positive probability satisfies
$f(P[W])>r/2$. Then there is a fixed induced $k$-connected subgraph
$H\subseteq P$ such that
\begin{equation}\label{small:region-expectation}
 \mathbb E f(H[W\cap V(H)])>r/4.
\end{equation}
The subgraph $H$ is chosen before any fresh copies of the experiment
are sampled.
\end{lemma}

\begin{proof}
Build a binary decomposition of $P$, making all choices before sampling.
A node is an induced subgraph. Stop at a node of order at most $2k$, and
also at a larger node which is $k$-connected. Otherwise choose a separator
of size less than $k$, partition the vertices outside it into two
nonempty anticomplete sides, and form two children by adding the
separator to each side.

There are $O(n^2+1)$ nodes. To see this, give a node of order $u>2k$
potential $u-k$. Each large child's potential is strictly smaller than
its parent's, since the other exclusive side is nonempty. If both
children are large, their potentials sum to at most their parent's:
their orders sum to $u+|S|<u+k$, where $S$ is the separator. If just one
child is large, its potential alone is at most the parent's. Thus at
each depth the sum of the large-node potentials is at most $n$, while
the depth of a large-node branch is at most $n$. The small leaves
contribute at most two children per large node, in addition to a possible
small root.

Every separator and every small leaf has order at most $2k$. Its expected
intersection size with $W$ is at most $2kp\le r/1000$.
Lemma~\ref{small:concentration} and a union bound over the decomposition
show that, with probability $1-o(1)$, all these intersections have fewer
than $r/100$ vertices. If every $k$-connected leaf had expected fractional
chromatic number at most $r/4$ after restriction to $W$, the same lemma
would give, again with probability $1-o(1)$, fractional chromatic number
less than $3r/10$ in every such leaf. A small leaf has this latter bound
as well, since its fractional chromatic number is at most its order.

Fix an outcome satisfying all these conclusions. Write $F=P[W]$.
The assumption of the lemma gives $f(F)>r/2$, whereas every leaf
restriction has fractional chromatic number less than $3r/10$.
We compare these bounds by replacing every vertex of $F$ with a clique
of the same large order. Ordinary-chromatic connectivity extraction in
this larger graph will produce a subgraph of large chromatic number that
cannot cross any of the blown-up separators. It must lie in one
blown-up leaf, giving the contradiction.

For a positive integer $m$, form $F^{(m)}$ by replacing each vertex of
$F$ with a clique of order $m$, and replacing each edge with a complete
join between its two replacement cliques. Every color class of $F^{(m)}$ projects to a stable
set of $F$ and uses at most one vertex of each replacement clique.
Dividing the multiplicities of these projected stable sets by $m$
therefore gives a fractional coloring of $F$. Consequently
\begin{equation}\label{small:blowup-lower}
 \chi(F^{(m)})\ge mf(F)>mr/2.
\end{equation}

For each leaf restriction, choose a fractional coloring of weight less
than $3r/10$. Round its finitely many stable-set weights up to multiples
of $1/m$. Give a distinct color slot to each unit of weight $1/m$.
Every original vertex belongs to at least $m$ slots, so its $m$ copies
can be assigned distinct slots containing that vertex. A slot is a
stable set upstairs: it contains at most one copy of any original
vertex, and its original vertices form a stable set. The rounding adds
only a bounded number of slots, independent of $m$ for this fixed leaf.
There are finitely many leaves, so for all sufficiently large $m$, every
blown-up leaf restriction has a coloring with fewer than $31mr/100$
colors.

Apply Theorem~\ref{thm:ordinary-extraction} to $F^{(m)}$ with the integer
parameter $k'=\lceil mr/50\rceil$. For sufficiently large $m$,
\eqref{small:blowup-lower} exceeds $7k'$. We obtain a $k'$-connected
subgraph $J$ with
\[
 \chi(J)\ge\chi(F^{(m)})-6k'
 >\frac{19mr}{50}-6
 >\frac{31mr}{100}.
\]
Each restricted separator of the decomposition has fewer than $r/100$
vertices, so its blow-up has fewer than $mr/100<k'$ vertices. Thus $J$
cannot meet both exclusive sides of any such split: deleting its
vertices in the separator would disconnect $J$ with fewer than $k'$
deletions. Nor can $J$ be contained in the separator, because its order
is greater than $k'$. Descending the decomposition places $J$ inside a
single blown-up leaf restriction, contradicting that leaf's coloring
bound.

It follows that some $k$-connected leaf satisfies
\eqref{small:region-expectation}. This argument also covers the case
where $P$ is too small to be $k$-connected: if all leaves were small, the
same contradiction would apply. Finally, the blow-up is only an
auxiliary graph for the extraction theorem, which requires no minor
exclusion; no bound on its clique-minor number was used.
\end{proof}

\subsection{Private models and a connected complement}

Assume the failure hypothesis of Lemma~\ref{small:fixed-region}, and fix
the resulting subgraph $H$. Sample $q$ fresh independent copies
$W_1,\ldots,W_q$ of \eqref{small:remainder}. Put
\[
 T_i=V(H)\cap\left(W_i\setminus\bigcup_{j\ne i}W_j\right),
 \qquad
 U=V(H)\cap\bigcup_{i=1}^qW_i.
\]
The sets $T_i$ are pairwise disjoint subsets of $U$.

\begin{lemma}\label{small:private}
With failure probability at most $4q(q-1)p+o(1)$, every $H[T_i]$ has
fractional chromatic number greater than $r/8$, and hence contains a
$K_{2x}$ model.
\end{lemma}

\begin{proof}
By \eqref{small:region-expectation} and
Lemma~\ref{small:concentration}, the events
\[
 \mathcal A_i=\{f(H[W_i\cap V(H)])>r/5\}
\]
all hold with probability $1-o(1)$: the individual failure probability
is at most $2\exp(-t^{7/10})$, and $q\le Q_0(\log t)^2$.

For each outcome of $W_i$ satisfying $\mathcal A_i$, choose optimal dual
weights on $H[W_i\cap V(H)]$ of total weight $a_i>r/5$.
Choose these weights as a function of $W_i$ alone. Conditional on this
one copy, each other copy still satisfies the vertex survival bound
\eqref{small:survival}. Therefore the expected weight on vertices
belonging to some other copy is at most $(q-1)pa_i$. Markov's inequality
shows that the probability of $\mathcal A_i$ together with the loss of
more than $a_i/4$ is at most $4(q-1)p$.

Taking a union bound over $i$, we conclude that, except with probability
$4q(q-1)p+o(1)$, all private sets retain weight at least
$3a_i/4>3r/20>r/8$. Restriction to $T_i$ preserves dual feasibility:
each stable set in $H[T_i]$ is also a stable set in
$H[W_i\cap V(H)]$. Hence $f(H[T_i])>r/8$. In this calculation we have
conditioned only on a single copy at a time, not on the simultaneous
events $\mathcal A_1,\ldots,\mathcal A_q$.

Theorem~\ref{thm:fractional} now gives
\[
 h(H[T_i])>r/32>2\lfloor r/200\rfloor=2x,
\]
so each private set contains the required model. The models are
vertex-disjoint because the private sets are disjoint.
\end{proof}

\begin{lemma}\label{small:complement}
With probability $1-o(1)$, every vertex of $H$ has at least $k/2$
neighbors in $H-U$, and $H-U$ is $\lfloor k/32\rfloor$-connected.
\end{lemma}

\begin{proof}
Each fixed vertex belongs to $U$ with probability at most
$qp\le10^{-4}$ by \eqref{small:sparseness}. Before sampling, fix $k$
neighbors of each vertex of $H$, possible since $H$ is $k$-connected.
The expected number of those neighbors in $U$ is at most $k/10000$.
Having at least $k/2$ of them in $U$ thus requires a deviation greater
than $k/3\ge r/3$. Lemma~\ref{small:concentration} and a union bound
over at most $n$ vertices prove the neighborhood assertion with
probability $1-o(1)$. This assertion holds for all vertices of $H$,
including those in $U$.

For the connectivity assertion, use the family of test matchings in
Lemma~\ref{lem:cut-test} for the fixed graph $H$, with order bound $n$
and connectivity parameter $k$. Fix this family before sampling the
fresh copies, and set
\[
 s=\left\lceil10\frac nk\log(2n)\right\rceil.
\]
The number of tests is at most
\begin{align*}
 (n+1)^{s+1}
 &=\exp\bigl(O((1+n/k)\log^2(2n))\bigr)\\
 &\le\exp\bigl(O(t^{1/10}\log^2t)\bigr).
\end{align*}
Here the constants are absolute, since $D_0$ is fixed.

The endpoint set of each test matching is a fixed set to which
Lemma~\ref{small:concentration} applies. It has size at most $k/4$, so
its expected intersection size with $U$ is at most $k/40000$. A count of
at least $k/32$ therefore requires a deviation greater than
$k/64\ge r/64$. Lemma~\ref{small:concentration} and the displayed bound
on the number of tests show that, with probability $1-o(1)$, every test
matching has fewer than $k/32$ endpoints in $U$. Together with the
neighborhood assertion, these are the hypotheses of
Lemma~\ref{lem:cut-test}, which gives the claimed connectivity.
\end{proof}

\begin{proof}[Proof of Lemma~\ref{small:eligible-thinning}]
If $P$ is empty, take all the stable sets empty. Otherwise use the
experiment \eqref{small:remainder} and suppose that every outcome of
positive probability has $f(P[W])>r/2$.
Lemma~\ref{small:fixed-region} supplies a fixed $H$. By
\eqref{small:sparseness}, the failure probability in
Lemma~\ref{small:private} is at most $4\cdot10^{-4}+o(1)$.
Thus Lemmas~\ref{small:private} and~\ref{small:complement} hold
simultaneously with positive probability for sufficiently large $t$.
Fix such an outcome, including its $q$ disjoint $K_{2x}$ models.

Partition the $t$ labels of a prospective $K_t$ model into
$b=\lceil t/x\rceil$ nonempty blocks of size at most $x$. Assign a
different private model to each single block and each unordered pair
of distinct blocks. There are
$b+\binom b2=b(b+1)/2\le b^2=q$ assignments. In a model assigned to
one block, retain one bag for each label of that block. In a model
assigned to two blocks, retain one bag for each label in their union.
At most $2x$ bags are retained from each model, and every pair of
distinct labels appears together in one assigned model.

Choose a port vertex in every retained bag. There are at most $2xq$
ports, and each has at least $k/2$ neighbors in $H-U$. By
\eqref{small:attachment-budget}, we can choose a distinct such neighbor
for every port, greedily. Let $m$ be the actual number of chosen
neighbors, so $1\le m\le2xq$. Their labels partition them into $t$
nonempty groups. The graph $H-U$ has connectivity at least
\[
 \left\lfloor\frac{k}{32}\right\rfloor
 \ge\gamma(2xq)\ge\gamma m,
\]
so Theorem~\ref{thm:grouping} provides pairwise disjoint connected sets
in $H-U$ containing these groups.

For each label, unite its connected set in $H-U$ with all its assigned
private bags and their port edges. This set is nonempty and connected.
Different labels give disjoint sets: their connected sets in $H-U$
are disjoint, all private bags are disjoint, and the bags lie in $U$.
Every two label sets are adjacent because their labels occur together
in an assigned clique model. These sets form a $K_t$ model in $P$,
contrary to $h(P)<t$. Some outcome of the sampling experiment must
therefore have fractional chromatic number at most $r/2$, as required.
\end{proof}

\subsection{Assigning the colors to bins}

\begin{proof}[Proof of Proposition~\ref{prop:small-list}]
Fix a sufficiently large integer $t$, and let $G$ satisfy the hypotheses
of the proposition. The empty graph needs no colors, so suppose that
$G$ is nonempty. Set
\[
 r_j=4t2^{-j},\qquad y_j=\max\{1,t/r_j\},
\]
and let $J$ be the largest nonnegative integer for which
$r_J\ge t/\log t$. For each $j\in\{0,\ldots,J\}$, allocate a separate
block of
\[
 N_j=\lceil Br_j\log(2y_j)\rceil
\]
bins. Allocate a further $16t$ bins for cleanup. All these bins are
distinct. Their total number $n_{\mathrm{bin}}$ is at most a fixed
constant times $t$. Indeed, $y_j\le2^j$, so
\begin{align*}
 n_{\mathrm{bin}}
 &\le4Bt\log2\sum_{j\ge0}(j+1)2^{-j}+(J+1)+16t\\
 &=16B(\log2)t+(J+1)+16t.
\end{align*}
Also $J+1=O(\log\log t)$, since $2^J\le4\log t$.
Choose an integer
\[
 D_{\mathrm{bin}}\ge16B\log2+18.
\]
For all sufficiently large $t$, we have
$n_{\mathrm{bin}}\le D_{\mathrm{bin}}t$.
Finally choose a positive integer $L$ such that
\begin{equation}\label{small:L-choice}
 \frac{L}{2D_{\mathrm{bin}}}>\log2+1.
\end{equation}

Let $\mathcal L$ be any finite list assignment on $G$ with
$|\mathcal L(v)|\ge Lt$ at every vertex. Map each color in the finite
union $\bigcup_v\mathcal L(v)$ independently and uniformly to one of
the $n_{\mathrm{bin}}$ bins. We claim that a map can be fixed such that
every vertex has an allowable color in at least half the bins of
every block, including cleanup.

For a fixed vertex and a block of $u$ bins, failure of this property
implies that a subset of $\lceil u/2\rceil$ bins is missed by its list.
For a specified such subset, the probability of being missed is at most
\[
 \left(1-\frac{u}{2n_{\mathrm{bin}}}\right)^{|\mathcal L(v)|}
 \le\exp\left(-\frac{Lt\,u}{2n_{\mathrm{bin}}}\right)
 \le\exp\left(-\frac{Lu}{2D_{\mathrm{bin}}}\right).
\]
There are at most $2^u$ choices for the missed subset, so
\eqref{small:L-choice} bounds the failure probability by $e^{-u}$.
Every nonfinal block has
$u\ge B(\log2)t/\log t$, and the cleanup block has $u=16t$.
Thus, with $c_B=\min\{B\log2,16\}>0$ and $\log t\ge1$, a union bound
over all vertices and blocks gives failure probability at most
\[
 (J+2)t^{11/10}\exp\left(-\frac{c_Bt}{\log t}\right)=o(1).
\]
For sufficiently large $t$, the desired map therefore exists; fix one.
This probability estimate depends on the list sizes, but not on the
number of distinct colors in their union.

Call a vertex eligible for a bin if its list has a color assigned to
that bin. It remains eligible whenever it survives to a later stage.
At the start, Theorem~\ref{thm:fractional} gives $f(G)\le4t=r_0$.
At stage $j$, let $P$ be the current induced remainder. If it is empty,
stop. Otherwise $h(P)<t$, $|V(P)|\le t^{11/10}$, and inductively
$f(P)\le r_j$. For each bin of that stage, take its eligible vertex set
inside $P$. Every vertex belongs to at least half these sets.
Lemma~\ref{small:eligible-thinning}, with $n=t^{11/10}$, supplies a
stable set for each bin such that deleting the eligible covered
vertices leaves fractional chromatic number at most
$r_j/2=r_{j+1}$. Retain one covering bin for every vertex removed at
this stage.

After stage $J$, the remaining graph has fractional chromatic number
at most $r_{J+1}<t/\log t$. If it is nonempty, put $\rho=t/\log t$ and
independently sample a stable set in each cleanup bin from a fractional
coloring distribution having every vertex inclusion probability at
least $1/\rho$. Such a distribution exists because $f\le\rho$.
Every remaining vertex is eligible in at least $8t$ cleanup bins, so
its probability of being missed by all eligible samples is at most
\[
 (1-1/\rho)^{8t}\le\exp(-8t/\rho)=t^{-8}.
\]
The probability that any remaining vertex is missed is therefore at
most $t^{11/10}t^{-8}=t^{-69/10}<1$. Choose an outcome covering every
remaining vertex and again retain one covering bin for each vertex.

Every vertex now has a retained bin and belongs to the stable set
associated with that bin. Choose any color in its original list which
maps to that bin. If two vertices receive the same color, the global
map sends that color to their common retained bin. Both vertices lie
in that bin's stable set and are therefore nonadjacent. This is a
proper $\mathcal L$-coloring of $G$.

The choices of constants were made in the order $\gamma$, $B$,
$D_{\mathrm{bin}}$, and $L$. All subsequent cutoffs depend only on these
constants, and every thinning estimate is uniform over the intermediate
graphs and eligibility sets. Enlarging one absolute $t_0$ to satisfy
the finitely many cutoff requirements proves the proposition.
\end{proof}
